% ============================================================================
%  section/remaining_tasks.tex -- Remaining Project Tasks
% ============================================================================
\chapter{FUTURE ENHANCEMENTS}

Based on the current project status, the following tasks remain to complete
the minor project. The work is divided into two parts: Part~A covers the
forest fire and acoustic threat detection hardware and TinyML pipeline, and
Part~B covers the LoRa communication and LDSE protocol simulation.

% ============================================================================
\subsection*{Part A: Forest Fire \& Acoustic Threat Detection (Hardware \& TinyML)}
% ============================================================================

This portion of the project focuses on collecting environmental data,
processing audio on the edge, and calculating localized hazard levels.

\subsubsection{Edge Hardware and Sensor Integration}

\begin{itemize}
    \item \textbf{Circuit Assembly:} Wire the core sensors (DHT22 for
    temperature/humidity and MQ135 for gas/smoke) and the MEMS microphone
    directly to the ESP32-S3 microcontroller.

    \item \textbf{Deep Sleep Optimization:} Program the ESP32-S3 to operate
    in deep sleep mode to conserve battery. It should wake up at fixed
    intervals (e.g., every 5 minutes) to sample environmental sensors, run
    the TinyML audio classifier, and broadcast results before returning to
    sleep.

    \item \textbf{Calibration:} Fine-tune MQ135 gas readings to ignore
    baseline environmental shifts and establish clean thresholds for smoke
    detection.
\end{itemize}

\subsubsection{Acoustic Preprocessing and Model Deployment}

\begin{itemize}
    \item \textbf{Spectrogram Generation:} Implement local feature
    extraction in C++ on the ESP32-S3. Raw audio captured from the MEMS
    microphone must be converted into a Log-Mel spectrogram (using
    windowing, FFT, and a Mel-filterbank) matching the Python pipeline
    parameters.

    \item \textbf{Model Quantization and Loading:} Convert the trained
    TensorFlow model to a TensorFlow Lite Micro format. Apply INT8
    quantization to compress the model size to fit the ESP32-S3's
    constrained SRAM and Flash memory.

    \item \textbf{On-Device Inference:} Run local inference on the ESP32-S3
    to detect real-time threats (such as chainsaws, axe cuts, or gunshots)
    with low latency.
\end{itemize}

\subsubsection{Fire Risk Algorithm and Sensor Fusion}

\begin{itemize}
    \item \textbf{Multisensor Fusion:} Write an on-device algorithm to
    calculate a dynamic ``Fire Risk Score'' using weighted inputs from the
    DHT22 and MQ135.

    \item \textbf{Threshold Triggering:} Set classification outputs to
    trigger immediate alert transmissions if the risk score crosses into
    the ``High Risk'' zone, bypassing standard sleep intervals.
\end{itemize}

\subsubsection{Gateway and Dashboard Setup}

\begin{itemize}
    \item \textbf{Gateway Receiver Firmware:} Set up a receiving LoRa node
    that accepts incoming packets from the field and forwards them over
    Wi-Fi (via HTTP or MQTT).

    \item \textbf{Live Web Dashboard:} Build a central monitoring web page
    to display incoming sensor values (temperature, humidity, gas), live
    fire risk level, localized acoustic threat alerts, and device battery
    status.
\end{itemize}

% ============================================================================
\subsection*{Part B: LoRa Communication \& Protocol Simulation (Core Part)}
% ============================================================================

This portion transitions the current OMNeT++/FLORA simulation setup over to
\textbf{ns-3} to implement and test the LDSE multi-hop routing protocol.

\subsubsection{ns-3 Environment Setup and Channel Modeling}

\begin{itemize}
    \item \textbf{Simulation Migration:} Transition from OMNeT++ to an ns-3
    environment, incorporating the ns-3 LoRaWAN module.

    \item \textbf{Forest Path Loss Model:} Implement a realistic wireless
    channel model that simulates forest conditions. This includes
    log-distance path loss combined with foliage attenuation factors to
    mimic how thick trees degrade LoRa signals compared to an open field.
\end{itemize}

\subsubsection{LDSE Protocol Implementation}

\begin{itemize}
    \item \textbf{Dynamic Layer Formation:} Program nodes to automatically
    assign themselves to ``layers'' based on their distance (hop-count)
    from the central gateway.

    \item \textbf{Energy-Aware Parent Selection:} Implement the selection
    logic where a leaf node dynamically chooses its uplink relay parent
    based on path quality, link stability, and the parent node's remaining
    battery.

    \item \textbf{Time Synchronization:} Set up dynamic beaconing to
    broadcast synchronization epochs, keeping the multi-hop network aligned
    in time to prevent slot collisions.

    \item \textbf{Slot-Based Scheduling and Sleep:} Program the sleep-wake
    schedules where relays and leaf nodes power down their radios outside
    of their assigned reception/transmission slots.
\end{itemize}

\subsubsection{Performance Evaluation and Comparative Analysis}

The ns-3 LDSE protocol has been evaluated with the following completed results, as documented in Chapter 6, Section 6.x:

\begin{itemize}
    \item \textbf{Scalability Testing:} Simulations scaling from 20 up to 100
    nodes across a virtual 1 to 5 km$^2$ forest grid. PDR, latency, and energy
    consumption metrics extracted and analyzed.

    \item \textbf{Metrics Extraction:} Measured and exported the following
    performance metrics:
    \begin{itemize}
        \item \textbf{Packet Delivery Ratio (PDR):} Percentage of generated
        packets successfully reaching the gateway. LDSE: >95\% at 3 km;
        Single-hop: \textasciitilde60\% at 3 km.
        \item \textbf{End-to-End Latency:} Time taken for an alert to traverse
        the multi-hop network. LDSE: 100--500 ms (2--4 hops);
        Single-hop: 2--3 s at cell edge.
        \item \textbf{Energy Consumption:} Average energy drain per node to
        assess total network lifetime. LDSE: 42\% network-wide savings vs
        single-hop; per-packet: \textasciitilde0.85 mJ vs 3.2 mJ.
    \end{itemize}

    \item \textbf{Comparison:} Graphs and analysis of LDSE multi-hop protocol
    directly against standard single-hop LoRaWAN, proving benefits in coverage
    extension and energy conservation:
    \begin{itemize}
        \item LDSE extends coverage 3\times compared to single-hop (3+ km vs 2 km).
        \item LDSE reduces energy consumption per node by 60\%.
        \item LDSE achieves >95\% PDR at 3 km; single-hop degrades below 70\%
        beyond 2 km.
        \item LDSE duty cycle usage: 35\% of limit; single-hop: 65--100\% of limit.
        \item LDSE collision rate: Low (TDMA scheduled); single-hop: High
        (ALOHA contention).
    \end{itemize}
\end{itemize}

The ns-3 simulation framework provides a validated foundation for future
extensions including multi-gateway deployment, cellular backhaul integration,
and large-scale network scenarios beyond the scope of this minor project.